% Generated draft: requires mathematical and source review.
\documentclass[11pt]{article}
\usepackage[margin=25mm]{geometry}
\usepackage{fontspec}
\setmainfont{Latin Modern Roman}
\usepackage{amsmath,amssymb,mathtools,unicode-math}
\setmathfont{Latin Modern Math}
\usepackage{enumitem,needspace,xurl,hyperref}
\setlength{\parindent}{0pt}
\setlength{\parskip}{5pt}
\setlength{\emergencystretch}{3em}
\newcommand{\N}{\mathbb{N}}
\newcommand{\Z}{\mathbb{Z}}
\newcommand{\Q}{\mathbb{Q}}
\newcommand{\R}{\mathbb{R}}
\newcommand{\C}{\mathbb{C}}
\newcommand{\F}{\mathbb{F}}
\newcommand{\A}{\mathbb{A}}
\newcommand{\PP}{\mathbb{P}}
\newcommand{\E}{\mathbb{E}}
\newcommand{\eps}{\varepsilon}
\newcommand{\dd}{\,\mathrm{d}}
\begin{document}
\section*{UnsolvedMath ID: 30001404}
\textbf{Source problem code:} catalogue entry
\subsection{Definitions and mathematical statement}
Let \(K\) be an infinite field equipped with additional structure such that its ambient first-order theory is stable and the structure is saturated, as stipulated in the supplied record. All references to genericity concern this ambient structure. The language, degree of saturation, precise meaning of ``small,'' and intended notion of genericity are not specified in the supplied excerpt.

Write \(\operatorname{Gen}_k(g)\) for the source's assertion that \(g\in K\) is generic over the small subfield \(k\subseteq K\). This notation records an unresolved source notion; it does not supply its missing definition. Likewise, let \(\mathsf{E}(K)\) denote the source's ``exchange property for generics,'' whose definition is absent from the supplied record.

On the field-theoretic reading of ``algebraic over \(k(h)\),'' define the second stated condition by
\[
\mathsf{P}(K):\quad
\forall\, k\subseteq K\text{ a small subfield}\;
\forall\,g,h\in K,\qquad
\bigl(\operatorname{Gen}_k(g)\ \land\ g\text{ is algebraic over }k(h)\bigr)
\Longrightarrow \operatorname{Gen}_k(h).
\]
Here \(k(h)\) is the subfield generated by \(k\cup\{h\}\), and field-theoretic algebraicity means that some nonzero polynomial \(f\in k(h)[X]\) satisfies \(f(g)=0\). Whether the source instead intends model-theoretic algebraic closure in the expanded language requires clarification.

The supplied corrected statement asks whether, for every such structured field \(K\), both implications hold:
\[
K\text{ separably closed}\ \Longrightarrow\ \mathsf{E}(K),
\qquad
K\text{ algebraically closed}\ \Longrightarrow\ \mathsf{P}(K).
\]
A field is separably closed if it has no proper algebraic separable extension, and algebraically closed if every nonconstant polynomial over it has a root in it.

There is an explicit directional conflict in the record: its original statement instead asks for
\[
\mathsf{E}(K)\ \Longrightarrow\ K\text{ separably closed},
\qquad
\mathsf{P}(K)\ \Longrightarrow\ K\text{ algebraically closed}.
\]
The supplied verification field describes these latter directions as proved in the report and the corrected directions as conjectural. This is a supplied source-verification claim, not an independently established fact here. The missing definitions and the referenced but unsupplied hypotheses of Conjectures 7 and 8 prevent a fully resolved formal statement.
\subsection{Short English statement}
The supplied corrected question asks whether separable closure forces exchange for generics, and whether algebraic closure forces an element to be generic whenever a generic element is algebraic over the field it generates with a small base subfield. The original statement reverses both implications. Definitions of genericity and exchange, together with the full hypotheses of the cited conjectures, are needed for mathematical review.
\subsection{Sources}
\begin{itemize}
\item Local UnsolvedMath registry, supplied record with permanent ID \(30001404\), exact source problem code \texttt{catalogue entry}, title ``Closure Implies Exchange for Generics,'' in the collection ``Oberwolfach Reports Open Problems.'' The record gives proposed year 2010 and names no individual proposer.
\item Supplied report citation: \emph{Model Theory: Around Valued Fields and Dependent Theories} (2010), Oberwolfach Reports, pp.\ 5--53, DOI \texttt{10.4171/owr/2010/01}. Source URL and supplied literature-source link, labeled ``original Oberwolfach problem report'': \url{https://doi.org/10.4171/owr/2010/01}. The background describes these reports as documenting mathematical workshops and their open-problem discussions.
\item Literature assessment dated exactly \texttt{2026-08-21}: the supplied curation labels the problem open and says that the extracted sentence leaves both converses open, with no proof of either converse in the surrounding report. This is supplied dated curation, not independently verified; the background expressly requests independent verification of time-sensitive status and cited claims.
\item The supplied statement status is \texttt{corrected\_verified}. Its verification text attributes proofs of the exchange-to-closure directions and conjectures in the reverse directions to the official report. These are supplied verification claims; the report and Conjectures 7 and 8 have not been independently inspected here.
\end{itemize}
\end{document}
